% Generated draft: requires mathematical and source review.
\documentclass[11pt]{article}
\usepackage[margin=25mm]{geometry}
\usepackage{fontspec}
\setmainfont{Latin Modern Roman}
\usepackage{amsmath,amssymb,mathtools,unicode-math}
\setmathfont{Latin Modern Math}
\usepackage{enumitem,needspace,xurl,hyperref}
\setlength{\parindent}{0pt}
\setlength{\parskip}{5pt}
\setlength{\emergencystretch}{3em}
\newcommand{\N}{\mathbb{N}}
\newcommand{\Z}{\mathbb{Z}}
\newcommand{\Q}{\mathbb{Q}}
\newcommand{\R}{\mathbb{R}}
\newcommand{\C}{\mathbb{C}}
\newcommand{\F}{\mathbb{F}}
\newcommand{\A}{\mathbb{A}}
\newcommand{\PP}{\mathbb{P}}
\newcommand{\E}{\mathbb{E}}
\newcommand{\eps}{\varepsilon}
\newcommand{\dd}{\,\mathrm{d}}
\begin{document}
\section*{UnsolvedMath ID: 30004493}
\textbf{Source problem code:} catalogue entry
\subsection{Definitions and mathematical statement}
Let \(\Phi:B\to\mathcal P\) denote the period map referred to in the supplied record, and write its stated Stein factorization as
\[
B\xrightarrow{\widehat\Phi}\widehat{\mathcal P}\xrightarrow{\nu}\mathcal P,
\qquad \Phi=\nu\circ\widehat\Phi.
\]
Here the Stein factorization terminology indicates that \(\widehat\Phi\) has connected fibers and \(\nu\) is finite. Let \(X=\widehat{\mathcal P}^{T}\) be the specified compact complex analytic variety containing \(\widehat{\mathcal P}\) as an open subspace. A compact complex analytic variety is projective if it admits a closed holomorphic embedding into complex projective space.

\textbf{Question.} Is this specified completion \(X\) projective; that is, does there exist an integer \(N\geq 0\) and a closed holomorphic embedding
\[
\iota:\widehat{\mathcal P}^{T}\hookrightarrow\PP^{N}(\C),
\]
so that \(\widehat{\mathcal P}^{T}\) is a projective completion of \(\widehat{\mathcal P}\)?

\textbf{Unresolved qualifications.} The source does not specify the base \(B\), the underlying variation of Hodge structure, the hypotheses on the period map, or the construction denoted by the superscript \(T\). The question therefore concerns the particular completion intended by the source, whose full definition is missing here, and does not quantify over arbitrary compact analytic spaces containing \(\widehat{\mathcal P}\). A completion ordinarily requires a dense open inclusion; the displayed source statement explicitly asserts only openness. It also does not specify whether this inclusion is analytically open or is intended to be Zariski open in a compatible algebraic structure. The displayed projectivity question expresses the stated central target, while these additional completion qualifications require clarification.
\subsection{Short English statement}
The question asks whether the particular compact analytic completion associated with the Stein factorization of a period map is projective. The supplied record leaves the period-map hypotheses and the construction of this completion unspecified. Its dated assessment reports partial progress, including projective fibers, while retaining projectivity of the completion as a conjecture; this assessment is not independently verified here.
\subsection{Sources}
\begin{itemize}
\item Supplied local UnsolvedMath registry record, permanent ID \(30004493\), title: \emph{Projectivity of the Stein-Factorized Period-Map Completion}; source problem code \texttt{catalogue entry}. The collection is \emph{Oberwolfach Reports Open Problems}. No individual proposer is identified; the supplied proposed year is 2020.
\item Supplied report citation: \emph{Komplexe Analysis -- Algebraicity and Transcendence} (2021), Oberwolfach Reports, pp.\ 1263--1313, DOI \texttt{10.4171/owr/2020/24}. \url{https://doi.org/10.4171/owr/2020/24}. The supplied background describes these reports as records of mathematical workshops and their open-problem discussions.
\item Supplied literature reference: \emph{Global asymptotic structure of period mappings}. \url{https://arxiv.org/abs/2010.06720}.
\item Supplied literature reference: \emph{Geometric and Hodge-theoretic completion of period mappings}. \url{https://arxiv.org/abs/2506.10109}.
\item Literature assessment dated exactly \texttt{2026-08-22}: the supplied curation labels the problem \texttt{partially\_solved} and describes the proposed completion as a Hausdorff topological and finite normal analytic completion with projective fibers, while reporting that its projective algebraicity remains conjectural in the primary literature. This is supplied dated curation derived from the OWR record, not an independently verified literature assessment. The supplied record also labels the statement \texttt{corrected\_verified} and reports prior manual semantic, typographic, and LaTeX review; those are registry annotations, not mathematical verification of this draft.
\end{itemize}
\end{document}
